% !TEX-PROGRAM: PDFLATEX

% DOCUMENT CLASS

\documentclass[12pt,letterpaper,twoside,final]{article}

% FONT ENCODING

\usepackage[T2C,T2B,T2A,LGR,TS1,T1]{fontenc}

% SÍMBOLOS

\usepackage[full]{textcomp}

% TYPEFACES

\usepackage[mainfont,semibold,scosf]{ysabeau}
\usepackage{hologo}

% LANGUAGE

\usepackage[american]{babel}

% PAGE LAYOUT

\usepackage[margin=2.5cm,includehead=true,headsep=1.5cm]{geometry}

% INDENT AND LINE SPACING

\usepackage[indent=5mm,skip=0pt]{parskip}

% MICROTYPOGRAPHY

\usepackage[activate={true,nocompatibility},
final=true,babel=true,tracking=true,kerning=true,spacing=true,
factor=1000,stretch=20,shrink=20]{microtype}
\SetTracking{encoding=*,shape=sc}{30}
%\DeclareMicrotypeSet*{smallcapsi}
%{encoding={TS1,T1},shape={sc*,si,scit}}

% FOOTNOTE MARKS

\let\oldfootnote\footnote
\renewcommand\footnote[1]{%
\oldfootnote{\hspace{.5mm}#1}}

% HEADERS

\usepackage[nocheck]{fancyhdr}

% TABLES

\usepackage{longtable,booktabs}
\usepackage[font=footnotesize]{caption}
\usepackage{float}

% LISTS

\usepackage{enumitem}

% FONT TABLES

\usepackage{fonttable}

% GRAPHICS

\usepackage{graphicx}

% HYPERLINKS

\usepackage[hyperindex=true,final=true,bookmarks=true,bookmarksnumbered=true,
bookmarksopen=true,breaklinks=true,citecolor=black,colorlinks=true,
linkcolor=black,urlcolor=black,pdftitle={The Ysabeau font package},pdfsubject={TeX},pdfauthor={Noel Merino Hernández},pdfkeywords={Typography Ysabeau Christian Thalmann LaTeX pdfLaTeX},pdfproducer={pdflatex 3.14159265-2.6-1.40.17 (TeX Live 2016/Debian)},pdfcreator={Noel Merino Hernández (muxkernel@gmail.com)}]{hyperref}

\urlstyle{sf}

\usepackage{hyperxmp}
\hypersetup{
pdfcopyright={This document is provided as is, but without warranty, and is distributed under Do What The Fuck You Want To Public License (WTFPL), version 2.0 or later.},
pdflicenseurl={http://www.wtfpl.net/}
}

\title{The Ysabeau font\footnote{version font: 2.002} package}
\author{Noel Merino Hernández \\ {\small \href{muxkernel@gmail.com}{muxkernel at gmail dot com}}}

% DOCUMENT
\begin{document}
\maketitle
\raggedbottom
\begin{abstract}
\noindent The Ysabeau font package provides \LaTeX\ support for traditional \TeX\ engines (\textsc{pdf}\TeX, dvips, an so on).\footnote{For shapes and weights, see \S~\ref{sec:shapes-and-weights}, page \pageref{sec:shapes-and-weights}.} For \hologo{XeTeX} and \hologo{LuaTeX} users, OpenType and TrueType fonts are \emph{only} provided to use with \texttt{fontspec} package. This package does not provide configuration files for those \TeX\ engines.
\end{abstract}
\section{Introduction}\label{sec:introduction}
\thispagestyle{empty}
\pagestyle{fancy}
\fancyhf{}
\fancyhead[CE,CO]{The Ysabeau font package}
\fancyhead[LE,RO]{\thepage}
\renewcommand{\headrulewidth}{0.4pt}
On github, we can read:
\begin{quote}
Ysabeau is a free type family developed by Christian Thalmann (Catharsis Fonts). It combines the time-honored and supremely readable letterforms of the Garamond legacy with the clean crispness of a low-contrast sans serif, rendering it well suited for body copy as well as display.
\end{quote}
Ysabeau font was drawn from scratch by Christian Thalmann, using the comercial font editor Glyphs, but was distributed under \textsc{sil} open font license. Previously known as ``Eua de Garamond'', the typeface was later renamed after Ysabeau, the second wife of Claude Garamond (1499--1561). In the end, it's just another Garamond typeface in sans serif form living out there.
\section{Document license}\label{sec:document-license}
This document is provided \emph{as is,} without warranty, and is distributed under \href{http://www.wtfpl.net}{Do What The Fuck You Want To Public License (\textsc{wtfpl})}, version 2 or later. For a detailed information about this license, visit \url{http://www.wtfpl.net/about}
\section{Interference with other packages}
The Ysabeau font package load automatically the following packages: \texttt{fontenc, textcomp, xstring, ifthen, scalefnt, etoolbox, xkeyval, mweights} and \texttt{fontaxes}. If you want to pass some options a these packages, just do it \emph{before} to use the \texttt{Ysabeau} font package. Example:
\begin{verbatim}
\usepackage[T2A,LGR,TS1,T1]{fontenc} % load before Ysabeau
\usepackage[full]{textcomp} % load before Ysabeau
\usepackage[osf,scosf]{ysabeau}
\end{verbatim}
\section{Usage}\label{sec:usage}
Just add \verb|\usepackage{ysabeau}| at the preamble of your document. That's it!
\section{Package options}\label{sec:package-options}
\begin{table}[H]
\centering
\caption{Package options}
\begin{tabular}{@{}lll@{}}
\toprule
Option & Example or description & Render as \\
\midrule
scale o scaled & \verb|\usepackage[scale=0.9]{ysabeau}| & {} \\
lining, lf or LF & it will use lining figures & \textlf{12345} \\
oldstyle, osf or OSF\textsu{*} & it will use oldstyle figures & \textosf{12345} \\
tabular, t or T & it will use tabular figures & \texttlf{12345} \\
proportional, p or P\textsu{*} & it will use proportional figures & {} \\
mainfont & it will use Ysabeau font (\verb|\rmfamily|) by default & {} \\
scosf & it will use oldstyle figures in small caps block & \textsc{xvii--xviii} \\
black & it will use black font (\verb|\bfseries|) & \black{black} \\
extrabold & it will use extrabold font (\verb|\bfseries|) & \extrabold{extrabold} \\
bold\textsu{*} & it will use bold font (\verb|\bfseries|) & \bold{bold} \\
semibold & it will use semibold font (\verb|\bfseries|) & \semibold{semibold} \\
medium\textsu{1} & it will use medium font (\verb|\mdseries|) & {} \\
regular\textsu{*} & it will use regular font (\verb|\mdseries|) & regular \\
light & it will use light font (\verb|\mdseries|) & \light{light} \\
extralight & it will use extralight font (\verb|\mdseries|) & \extralight{extralight} \\
thin & it will use thin font (\verb|\mdseries|) & \thin{thin} \\
hairline & it will use hairline font (\verb|\mdseries|) & \hairline{hairline} \\
\bottomrule
\end{tabular}
\caption*{* default.\hspace*{1em}\textsu{1} a bit heavier than regular.}
\end{table}
\section{Font encodings}\label{sec:font-encodings}
\begin{table}[H]
\centering
\caption{Font encodings}
\begin{tabular}{@{}ll@{}}
\toprule
Language & Encodings \\
\midrule
Vietnamese & \textlf{T5} \\
Afrikáans & \texttlf{T4} \\
\textsc{ipa} & \texttlf{TS3, T3} \\
Latin & \textlf{QX, CS, L7X, LY1, OT1, TS1, T1} \\
Cyrillic & \texttlf{T2C, T2B, T2A} \\
Greek & \texttlf{LGR} \\
\bottomrule
\end{tabular}
\end{table}
\section{Shapes and weights}\label{sec:shapes-and-weights}
\begin{table}[H]
\centering
\caption{Shapes}
\begin{tabular}{@{}ll@{}}
\toprule
Shape & Render as \\
\midrule
n & Romans \\
it & \emph{Italics} \\
sw & \textsw{ct, fi, fl, rt, sp, \&, Q} \\
%sw & \textsw{Calligraphics} \\
sc & \textsc{Small caps} \\
\bottomrule
\end{tabular}
\label{tab:shapes}
\end{table}
\begin{table}[H]
\centering
\caption{Weights}
\begin{tabular}{@{}ll@{}}
\toprule
Weight & Render as \\
\midrule
Hairline & \hairline{Hairline} \\
Thin & \thin{Thin} \\
Extralight & \extralight{Extralight} \\
Light & \light{Light} \\
Regular & Regular \\
Medium\textsu{1} & {} \\
Semibold & \semibold{Semibold} \\
Bold & \bold{Bold} \\
Extrabold & \extrabold{Extrabold} \\
Black & \black{Black} \\
\bottomrule
\end{tabular}
\caption*{\textsu{1} a bit heavier than regular.}
\label{tab:weights}
\end{table}
\noindent The following commands and declarations (Table~\ref{tab:commands-for-weights}) are provided \emph{just} for local use (weights), regardless of the options chosen in the package.
\begin{table}[H]
\centering
\caption{Commands and declarations for weights}
\begin{tabular}{@{}llll@{}}
\toprule
Weight & Command & Declaration & Render as \\
\midrule
Hairline & \verb|\hairline{}| & \verb|\hairlinefont | & {\hairlinefont Hairline} \\
Thin & \verb|\thin{}| & \verb|\thinfont | & {\thinfont Thin} \\
Extralight & \verb|\extralight{}| & \verb|\extralightfont | & {\extralightfont Extralight} \\
Light & \verb|\light{}| & \verb|\lightfont | & {\lightfont light} \\
Semibold & \verb|\semibold{}| & \verb|\semiboldfont | & {\semiboldfont Semibold} \\
Bold & \verb|\bold{}| & \verb|\boldfont | & {\boldfont Bold} \\
Extrabold & \verb|\extrabold{}| & \verb|\extraboldfont | & {\extraboldfont Extrabold} \\
Black & \verb|\black{}| & \verb|\blackfont | & {\blackfont Black} \\
\bottomrule
\end{tabular}
\label{tab:commands-for-weights}
\end{table}
\section{Figures}\label{sec:figures}
See also package options (\S\,\ref{sec:package-options}, page~\pageref{sec:package-options}) for figures. The following commands and declarations (Table \ref{tab:figures}) are provided \emph{just} for local use, regardless of the options chosen in the package.
\begin{table}[h]
\centering
\caption{Figures}
\begin{tabular}{@{}llr@{}}
\toprule
Command & Declaration & Render as \\
\midrule
\verb|\textlf{}| & \verb|\lfstyle| & {\lfstyle 1234567890} \\
\verb|\liningnums{}| & \verb|\lfstyle| & \liningnums{1234567890} \\
\verb|\texttlf{}| & \verb|\tlfstyle| & {\tlfstyle 1234567890} \\
\verb|\tabularliningnums{}| & \verb|\tlfstyle| & \tabularliningnums{1234567890} \\
\verb|\textosf{}| & \verb|\osfstyle| & {\osfstyle 1234567890} \\
\verb|\oldstylenums{}| & \verb|\osfstyle| & \oldstylenums{1234567890} \\
\verb|\texttosf{}| & \verb|\tosfstyle| & {\tosfstyle 1234567890} \\
\verb|\tabularoldstylenums{}| & \verb|\tosfstyle| & \tabularoldstylenums{1234567890} \\
\bottomrule
\end{tabular}
\label{tab:figures}
\end{table}
\subsection{Superiors}\label{subsec:superiors}
\begin{table}[H]
\centering
\caption{Superiors}
\begin{tabular}{@{}llr@{}}
\toprule
Command & Declarations & Render as \\
\midrule
\verb|\textsu{}| & \verb|\supfigures| or \verb|\supstyle| & \textsu{1234567890} \\
\verb|\textsup{}| & \verb|\supfigures| or \verb|\supstyle| & \textsup{1234567890} \\
\verb|\textsuperior{}| & \verb|\supfigures| or \verb|\supstyle| & \textsuperior{1234567890} \\
\verb|\supscript{}| & \verb|\supfigures| or \verb|\supstyle| & \supscript{1234567890} \\
\verb|\textsuperscript{}| & {} & \textsuperscript{1234567890} \\
\verb|\textsuperscript*{}| & {} & \textsuperscript*{1234567890} \\



\bottomrule
\end{tabular}
\end{table}
\subsection{Numerators}\label{subsec:numerators}
\begin{table}[H]
\centering
\caption{Numerators}
\begin{tabular}{@{}llr@{}}
\toprule
Command & Declarations & Render as \\
\midrule
\verb|\textnum{}| & \verb|\numfigures| or \verb|\numstyle| & \textnum{1234567890} \\
\verb|\textnumerators{}| & \verb|\numfigures| or \verb|\numstyle| & {\numstyle 1234567890} \\
\bottomrule
\end{tabular}
\caption*{\textsc{Note:} superiors and numerators are different: a\textsu{1}\textnum{1}b.}
\end{table}
\subsection{Denominators}\label{subsec:denominators}
\begin{table}[H]
\centering
\caption{Denominators}
\begin{tabular}{@{}llr@{}}
\toprule
Command & Declarations & Render as \\
\midrule
\verb|\textde{}| & \verb|\denfigures| or \verb|\denstyle| & \textde{1234567890} \\
\verb|\textden{}| & \verb|\denfigures| or \verb|\denstyle| & \textden{1234567890} \\
\verb|\textdenominators{}| & \verb|\denfigures| or \verb|\denstyle| & \textdenominators{1234567890} \\
\bottomrule
\end{tabular}
\end{table}
\subsection{Inferiors}\label{subsec:inferiors}
\begin{table}[H]
\centering
\caption{Inferiors}
\begin{tabular}{@{}llr@{}}
\toprule
Command & Declarations & Render as \\
\midrule
\verb|\textinf{}| & \verb|\inffigures| or \verb|\infstyle| & \textinf{1234567890} \\
\verb|\textinferior{}| & \verb|\inffigures| or \verb|\infstyle| & \textinferior{1234567890} \\
\verb|\subscript{}| & \verb|\inffigures| or \verb|\infstyle| & \subscript{1234567890} \\
\verb|\textsubscript{}| & {} & \textsubscript{1234567890} \\
\verb|\textsubscript*{}| & {} & \textsubscript*{1234567890} \\
\bottomrule
\end{tabular}
\caption*{\textsc{Note:} denominators and inferiors are different: a\textde{1}\textinf{1}b.}
\end{table}
\def\solidus{{\fontencoding{TS1}\selectfont\char47}}
\section{\texorpdfstring{\protect\textsuperscript*{Text}}{Text}\texorpdfstring{\protect\solidus}{\slash}\texorpdfstring{\protect\textsubscript*{fractions}}{fractions}}
\label{sec:fractions}
\begin{table}[H]
\centering
\caption{\protect\textsuperscript*{Text}\solidus\protect\textsubscript*{fractions}}
\begin{tabular}{@{}ll@{}}
\toprule
Command or declaration & Render as \\
\midrule
\verb|\textfrac[1]{1}{2}| & \textfrac[1]{1}{2} \\
\verb|\textfrac{1}{2}| & \phantom{0}\textfrac{1}{2} \\
\verb|\onequarter| & \phantom{0}\onequarter \\
\verb|\onehalf| & \phantom{0}\onehalf \\
\verb|\threequarters| & \phantom{0}\threequarters \\
\bottomrule
\end{tabular}
\end{table}
\section{Concluding remarks}\label{sec:concluding-remarks}
To use Ysabeau fonts, just add the line
\begin{verbatim}
\usepackage{ysabeau}
\end{verbatim}
at the preamble of your document. That's it.

The Ysabeau fonts were created to Adobe PostScript (type1) using \texttt{cfftot1}; \LaTeX\ files support were created with \texttt{autoinst}, wrapper of \texttt{LCDF TypeTools} suite. This package would not exist without the help of \TeX\ h4ck3rs: Marc Penninga (\texttt{autoinst}) and Eddie Kohler (\texttt{LCDF TypeTools}). Thank you so much guys.

The Ysabeau fonts are distributed under \textls[50]{SIL} open font license, version 1.1. A verbatim of this license can be found in \texttt{doc} folder. The \LaTeX\ files support are Public Domain. This document is licensed under \href{http://www.wtfpl.net}{Do What The Fuck You Want To Public License (\textsc{wtfpl})}, version 2 or later. A verbatim of this license can be found in \texttt{doc} folder. For a detailed information about this license, visit \url{http://www.wtfpl.net/about}
\section{Glyphs}\label{sec:glyphs}
\subsection{QX}
\fonttable{Ysabeau-Regular-lf-qx.tfm}
\subsection{\textlf{L7X}}
\fonttable{Ysabeau-Regular-lf-l7x.tfm}
\subsection{CS}
\fonttable{Ysabeau-Regular-lf-cs.tfm}
\subsection{\textlf{LGR}}
\fonttable{Ysabeau-Regular-lf-lgr.tfm}
\subsection{\textlf{T5}}
\fonttable{Ysabeau-Regular-lf-t5.tfm}
\subsection{\textlf{T4}}
\fonttable{Ysabeau-Regular-lf-t4.tfm}
\subsection{\textlf{TS3}}
\fonttable{Ysabeau-Regular-lf-ts3.tfm}
\subsection{\textlf{T3}}
\fonttable{Ysabeau-Regular-lf-t3.tfm}
\subsection{\textlf{T2C}}
\fonttable{Ysabeau-Regular-lf-t2c.tfm}
\subsection{\textlf{T2B}}
\fonttable{Ysabeau-Regular-lf-t2b.tfm}
\subsection{\textlf{T2A}}
\fonttable{Ysabeau-Regular-lf-t2a.tfm}
\subsection{\textlf{LY1}}
\fonttable{Ysabeau-Regular-lf-ly1.tfm}
\subsection{\textlf{OT1}}
\fonttable{Ysabeau-Regular-lf-ot1.tfm}
\subsection{\textlf{TS1}}
\fonttable{Ysabeau-Regular-lf-ts1.tfm}
\subsection{\textlf{T1}}
\fonttable{Ysabeau-Regular-lf-t1.tfm}
\subsection{\textlf{T1} (Swash)}
\fonttable{Ysabeau-Italic-lf-swash-t1.tfm}
\end{document}
